% Formulario T-17. Protocolo de aceptación (FEP01 p. 64).
\begin{formulario}{T-17}
\begin{contenidoExigido}{T-17}
  \exigencia{Entregables de cada hito y del producto final.}{La secuencia de actividades, dependencias y evidencias por hito se describe en \verSeccionFolio{sec:9-hitos}. La identificación final de entregables se concilia con la línea base contractual y el cronograma aprobado.}
  \exigencia{Criterios de aceptación objetivos.}{Se aplican la disponibilidad mínima de 99,9\,\% (\art{20}{14}), RTO máximo de 4 h y RPO máximo de 15 min (\caso{RT-07.04}{17}), y retención local mínima de 72 h (\caso{RT-03.10}{31}); véase también \verSeccionFolio{sec:9-metricas} y \verSeccionFolio{sec:9-criterios}. Los demás umbrales se aprueban y documentan antes de ejecutar las pruebas.}
  \exigencia{Evidencia requerida.}{Cada evidencia registra requisito y caso, fecha, versión, ambiente, datos y parámetros, herramienta, resultado, defectos y responsable; véase \verSeccionFolio{sec:9-evidencias}.}
  \exigencia{Plazos de revisión.}{Se aplican los plazos de revisión y subsanación establecidos en las bases y el contrato. Las ventanas y fechas por hito se incorporan al confirmarse la Carta Gantt; véase \verSeccionFolio{sec:9-hitos}.}
  \exigencia{Procedimiento de observaciones.}{Las observaciones se registran con identificador, requisito relacionado, severidad, evidencia y respuesta; la subsanación se revisa contra el criterio aprobado y conforme al procedimiento contractual; véase \verSeccionFolio{sec:9-protocolo-aceptacion}.}
  \exigencia{Acta de conformidad.}{El cierre registra el hito y versión evaluados, casos ejecutados y su estado, defectos abiertos, métricas efectivamente medidas, evidencias y pronunciamiento de la contraparte; véase \verSeccionFolio{sec:9-protocolo-aceptacion}.}
  \exigencia{Estado del catálogo de pruebas.}{El borrador contiene 120 casos inventariados: 25 unitarios, 25 de integración, 20 de sistema, 20 no funcionales y de seguridad, 15 HIL y 15 UAT. No se declaran ejecutados ni aprobados. Los escenarios sintéticos y umbrales no contractuales requieren fuente o aprobación antes de usarse como criterios de aceptación.}
\end{contenidoExigido}
\end{formulario}
